\chapter{Complexity — Machine}
%%% generated by the blueprint pipeline from TCSlib.Complexity.SpaceComplexity.Codec.Machine
\section{Overview}

Declarations of the Lean module \texttt{TCSlib.Complexity.SpaceComplexity.Codec.Machine}.

\section{Declarations}

\begin{definition}[First index in a two-element set]
\lean{Complexity.Codec.jLen}
\label{Complexity.Codec.jLen}
\leanok
Let $j$ be the first element of the two-element index set $	ext{Fin}(2)$; thus $j=0$.
\end{definition}

\begin{definition}[Second index in a two-element set]
\lean{Complexity.Codec.jBit}
\label{Complexity.Codec.jBit}
\leanok
Let $j$ be the second element of the two-element index set $	ext{Fin}(2)$; thus $j=1$.
\end{definition}

\begin{definition}[Code modes]
\lean{Complexity.Codec.CMode}
\label{Complexity.Codec.CMode}
\leanok
Define a four-element type of code modes, with one mode for valid codes, one for the
initial code, one for accepting codes, and one for step pairs.
\end{definition}

\begin{definition}[Code-part labels]
\lean{Complexity.Codec.CpId}
\label{Complexity.Codec.CpId}
\leanok
Let $
\mathcal{C}$ be the finite type whose 26 distinct elements are the labels “A,” “B,”
“Q1s,” “Cap,” “Q2,” “N1,” “Tn,” “E,” “Ts,” “F,” “Tw,” “U,” “Q1f,” “E2,” “Ts2,” “F2,”
“Tw2,” “U2,” “Q1f2,” “Tpos,” “Q3,” “Q4,” “Q1a,” “Q2a,” “Uc,” and “Tsc.”
\end{definition}

\begin{definition}[Six code labels]
\lean{Complexity.Codec.ScId}
\label{Complexity.Codec.ScId}
\leanok
Let $\mathrm{ScId}$ be the finite type consisting of six distinct labels, denoted $E$,
$F$, $E_2$, $F_2$, $\mathrm{In}$, and $C$.
\end{definition}

\begin{definition}[Instruction code labels]
\lean{Complexity.Codec.IcId}
\label{Complexity.Codec.IcId}
\leanok
Let $\mathrm{IcId}$ be the finite type with five constructors, denoted $E$, $U$, $E_2$,
$U_2$, and $C$.
\end{definition}

\begin{definition}[Data codec labels]
\lean{Complexity.Codec.DcId}
\label{Complexity.Codec.DcId}
\leanok
Let $\mathrm{DcId}$ be the three-element type whose elements are $S$, $Pos$, and $Pos2$.
\end{definition}

\begin{definition}[Comparison phase tags]
\lean{Complexity.Codec.PmPh}
\label{Complexity.Codec.PmPh}
\leanok
The type of comparison-phase tags has eight forms: `start`, `eq`, `inc`, `eq2`, `decR`
and `incR` each carrying a Boolean parameter, and `decN` and `eq3`.
\end{definition}

\begin{definition}[Virtual-adder states]
\lean{Complexity.Codec.VAd}
\label{Complexity.Codec.VAd}
\leanok
\uses{Complexity.Codec.AdSt, Complexity.Codec.VSt, Complexity.Codec.hH}
For a Boolean nondeterministic Turing machine $N$, a virtual-adder state is an augmented
state whose underlying state is a virtual-machine state, with a pointer and an array of
$h_H(N)+3k$ cells; the pointer ranges from $0$ through $h_H(N)+3k$, and each cell stores
a pair of Booleans.
\end{definition}

\begin{definition}[Simulation-adder states]
\lean{Complexity.Codec.SAd}
\label{Complexity.Codec.SAd}
\leanok
\uses{Complexity.Codec.AdSt, Complexity.Codec.SSt, Complexity.Codec.hH}
For a Boolean nondeterministic Turing machine $N$, a simulation-adder state is an
augmented state whose underlying state is a simulation state, with a pointer and an
array of $h_H(N)+3k$ cells; the pointer ranges from $0$ through $h_H(N)+3k$, and each
cell stores a pair of Booleans.
\end{definition}

\begin{definition}[Machine simulation labels]
\lean{Complexity.Codec.Lab}
\label{Complexity.Codec.Lab}
\leanok
\uses{Complexity.Codec.CpId, Complexity.Codec.DcId, Complexity.Codec.IcId, Complexity.Codec.PmPh, Complexity.Codec.SAd, Complexity.Codec.SInfo, Complexity.Codec.ScId, Complexity.Codec.VAd, Complexity.Codec.VSt, Complexity.Codec.W}
For a Boolean nondeterministic Turing machine $N$, the label type is generated by the
rejection and acceptance labels; labels for length scanning and incrementing; and
parameterized labels for the prefix, copying, counting, decrementing, simulation, and
input-processing routines. Its constructors also include labels indexed by states or
configurations of $N$, finite register indices, Boolean flags, and the comparison-phase
tags.
\end{definition}

\begin{definition}[Unconditional label jump]
\lean{Complexity.Codec.goto}
\label{Complexity.Codec.goto}
\leanok
\uses{Complexity.Codec.Lab, Complexity.Codec.rZ, Complexity.LogProg.Ins}
For a Boolean nondeterministic Turing machine $N$ and a label $\ell$, the instruction
associated with $\ell$ is a zero-test on the distinguished register $r_Z$ whose two
branches both go to $\ell$.
\end{definition}

\begin{definition}[Copy destination register]
\lean{Complexity.Codec.cpDst}
\label{Complexity.Codec.cpDst}
\leanok
\uses{Complexity.Codec.CpId, Complexity.Codec.rA, Complexity.Codec.rB, Complexity.Codec.rCap, Complexity.Codec.rE, Complexity.Codec.rE2, Complexity.Codec.rF, Complexity.Codec.rF2, Complexity.Codec.rN1, Complexity.Codec.rQ1, Complexity.Codec.rQ2, Complexity.Codec.rQ3, Complexity.Codec.rQ4, Complexity.Codec.rT, Complexity.Codec.rU}
For a Boolean nondeterministic Turing machine $N$, the copy-identifier type has the
constructors
$A,B,Q1s,Cap,Q2,N1,Tn,E,Ts,F,Tw,U,Q1f,E2,Ts2,F2,Tw2,U2,Q1f2,Tpos,Q3,Q4,Q1a,Q2a,Uc,Tsc$.
The destination-register map sends these, respectively, to
$r_A,r_B,r_{Q1},r_{Cap},r_{Q2},r_{N1},r_T,r_E,r_T,r_F,r_T,r_U,r_{Q1},r_{E2},r_T,r_{F2},r_T,r_U,r_{Q1},r_T,r_{Q3},r_{Q4},r_{Q1},r_{Q2},r_U,r_T$.
\end{definition}

\begin{definition}[Copy source register]
\lean{Complexity.Codec.cpSrc}
\label{Complexity.Codec.cpSrc}
\leanok
\uses{Complexity.Codec.CpId, Complexity.Codec.rA, Complexity.Codec.rB, Complexity.Codec.rE, Complexity.Codec.rE2, Complexity.Codec.rF, Complexity.Codec.rF2, Complexity.Codec.rL, Complexity.Codec.rN1, Complexity.Codec.rP, Complexity.Codec.rPos, Complexity.Codec.rS, Complexity.Codec.rWp}
For a Boolean nondeterministic Turing machine $N$, the copy-identifier type has the
constructors
$A,B,Q1s,Cap,Q2,N1,Tn,E,Ts,F,Tw,U,Q1f,E2,Ts2,F2,Tw2,U2,Q1f2,Tpos,Q3,Q4,Q1a,Q2a,Uc,Tsc$.
The source-register map sends these, respectively, to
$r_P,r_P,r_P,r_L,r_L,r_A,r_{N1},r_A,r_S,r_E,r_{Wp},r_F,r_F,r_B,r_S,r_{E2},r_{Wp},r_{F2},r_{F2},r_{Pos},r_A,r_B,r_A,r_A,r_A,r_S$.
\end{definition}

\begin{definition}[Initial virtual-adder configuration]
\lean{Complexity.Codec.vAd0}
\label{Complexity.Codec.vAd0}
\leanok
\uses{Complexity.Codec.VAd, Complexity.Codec.vInit}
For a Boolean nondeterministic Turing machine $N$, the initial virtual-adder
configuration has underlying state $vInit(N)$, pointer $0$, and every cell stores
$(\mathrm{false},\mathrm{false})$.
\end{definition}

\begin{definition}[Initial address state]
\lean{Complexity.Codec.sAd0}
\label{Complexity.Codec.sAd0}
\leanok
\uses{Complexity.Codec.SAd, Complexity.Codec.sInit}
For a nondeterministic Turing machine $N$ with Boolean tape alphabet, the initial
address state consists of two copies of the initial state $sInit(N)$, head position $0$,
and a tape of length $hH(N)+3k$ whose entries are all $(\mathrm{false},\mathrm{false})$.
\end{definition}

\begin{definition}[Successor labels for copy identifiers]
\lean{Complexity.Codec.cpNext}
\label{Complexity.Codec.cpNext}
\leanok
\uses{Complexity.Codec.CpId, Complexity.Codec.Lab, Complexity.Codec.W, Complexity.Codec.vAd0}
Let $N$ be a nondeterministic Turing machine with Boolean tape alphabet. Define a
function from the copy identifiers to control labels by assigning
$A,B,Q1s,Cap,Q2,N1,Tn,E,Ts,F,Tw,U,Q1f$ respectively to the labels for scanning the first
and second registers, copying the capacity, testing the counter, testing the $S$
down-counter, halving $N$, testing the halving counter, copying $Ts$, scanning $E$,
copying $Tw$, scanning $F$, incrementing $U$ at position $0$, and testing the $Pos$
down-counter. Assign $E2,Ts2,F2,Tw2,U2,Q1f2,Tpos,Q3,Q4,Q1a,Q2a,Uc,Tsc$ respectively to
copying $Ts2$, scanning $E2$, copying $Tw2$, scanning $F2$, incrementing $U2$ at
position $0$, testing the $Pos2$ down-counter, decrementing the input counter, copying
$Q4$, clearing the first verification register, copying $Q2a$, testing the initial
verification state, copying $Tsc$, and scanning $C$.
\end{definition}

\begin{definition}[Decrement-control pointer]
\lean{Complexity.Codec.dcPtr}
\label{Complexity.Codec.dcPtr}
\leanok
\uses{Complexity.Codec.DcId, Complexity.Codec.rQ1, Complexity.Codec.rQ2}
For a nondeterministic Turing machine $N$ with Boolean tape alphabet, the
decrement-control pointer sends $S$ to the register position $16$ and sends both $Pos$
and $Pos2$ to position $15$, among the $21$ register positions.
\end{definition}

\begin{definition}[Decrement-control base positions]
\lean{Complexity.Codec.dcBase}
\label{Complexity.Codec.dcBase}
\leanok
\uses{Complexity.Codec.DcId, Complexity.Codec.rE, Complexity.Codec.rE2, Complexity.Codec.rP}
For a nondeterministic Turing machine $N$ with Boolean tape alphabet, the base-position
map sends $S$ to register position $1$, $Pos$ to position $7$, and $Pos2$ to position
$9$, among the $21$ register positions.
\end{definition}

\begin{definition}[Decrement-control accumulator positions]
\lean{Complexity.Codec.dcAcc}
\label{Complexity.Codec.dcAcc}
\leanok
\uses{Complexity.Codec.DcId, Complexity.Codec.rPos, Complexity.Codec.rPos2, Complexity.Codec.rS}
For a nondeterministic Turing machine $N$ with Boolean tape alphabet, the
accumulator-position map sends $S$ to register position $4$, $Pos$ to position $11$, and
$Pos2$ to position $12$, among the $21$ register positions.
\end{definition}

\begin{definition}[Decrement-control strides]
\lean{Complexity.Codec.dcStride}
\label{Complexity.Codec.dcStride}
\leanok
\uses{Complexity.Codec.DcId}
For a nondeterministic Turing machine $N$ with Boolean tape alphabet, the stride
assigned to $S$ is $1$, while the strides assigned to $Pos$ and $Pos2$ are both $2$.
\end{definition}

\begin{lemma}[Upper bound on decrement strides]
\lean{Complexity.Codec.dcStride_le}
\label{Complexity.Codec.dcStride_le}
\leanok
\uses{Complexity.Codec.DcId, Complexity.Codec.dcStride}
\difficulty{4}
For every decrement instruction identifier, its stride is at most $2$.
\end{lemma}

\begin{definition}[Decrement-control descent labels]
\lean{Complexity.Codec.dcDowns}
\label{Complexity.Codec.dcDowns}
\leanok
\uses{Complexity.Codec.DcId, Complexity.Codec.Lab, Complexity.Codec.dcStride, Complexity.Codec.dcStride_le}
For a nondeterministic Turing machine $N$ with Boolean tape alphabet, a
decrement-control identifier $c$ and natural number $i$ determine a label as follows: if
$i$ is less than the stride assigned to $c$, the label is the descent label for $c$ at
index $i$; otherwise, it is the read label for $c$.
\end{definition}

\begin{definition}[Decrement-control increment labels]
\lean{Complexity.Codec.dcIncs}
\label{Complexity.Codec.dcIncs}
\leanok
\uses{Complexity.Codec.DcId, Complexity.Codec.Lab}
For a nondeterministic Turing machine $N$ with Boolean tape alphabet, a
decrement-control identifier $c$, a Boolean value $b$, and a natural number $i$
determine a label as follows: if $i<2$, the label is the increment label for $c$ and $b$
at index $i$; otherwise, it is the scale label for $c$ and $b$.
\end{definition}

\begin{definition}[Next labels for counter operations]
\lean{Complexity.Codec.dcNext}
\label{Complexity.Codec.dcNext}
\leanok
\uses{Complexity.Codec.DcId, Complexity.Codec.Lab}
For a nondeterministic Turing machine $N$ with Boolean tape alphabet, define a
label-valued map on the three counter identifiers by sending $S$ to the copy label for
$N1$, $Pos$ to the first length-test label, and $Pos2$ to the input-zero label.
\end{definition}

\begin{definition}[Source register for scaling operations]
\lean{Complexity.Codec.scSrc}
\label{Complexity.Codec.scSrc}
\leanok
\uses{Complexity.Codec.ScId, Complexity.Codec.rT}
For a nondeterministic Turing machine $N$ with Boolean tape alphabet, define the
source-register map for scaling identifiers to be constant: every identifier is assigned
register $13$ among the $21$ registers.
\end{definition}

\begin{definition}[Destination registers for scaling operations]
\lean{Complexity.Codec.scDst}
\label{Complexity.Codec.scDst}
\leanok
\uses{Complexity.Codec.ScId, Complexity.Codec.rE, Complexity.Codec.rE2, Complexity.Codec.rF, Complexity.Codec.rF2, Complexity.Codec.rU}
For a nondeterministic Turing machine $N$ with Boolean tape alphabet, define the
destination-register map on the six scaling identifiers by assigning $E,F,E2,F2,In,C$
respectively to registers $7,8,9,10,14,14$ among the $21$ registers.
\end{definition}

\begin{definition}[Scaling factors]
\lean{Complexity.Codec.scFactor}
\label{Complexity.Codec.scFactor}
\leanok
\uses{Complexity.Codec.ScId, Complexity.Codec.W}
For a nondeterministic Turing machine $N$ with Boolean tape alphabet, let $h$ be the
number of pairs consisting of an optional state of $N$ and one of three head positions,
and put $W=h+3k+1$, where $k$ is the number of tapes. Define the scaling factors for
$E,F,E2,F2,In,C$ to be respectively $4W,2,4W,2,2,2W$.
\end{definition}

\begin{lemma}[Upper bound for scan factors]
\lean{Complexity.Codec.scFactor_le}
\label{Complexity.Codec.scFactor_le}
\leanok
\uses{Complexity.Codec.ScId, Complexity.Codec.W, Complexity.Codec.scFactor}
\difficulty{4}
For every finite nondeterministic Turing machine $N$ and every scan identifier $c$, the
scan factor associated with $c$ is at most four times the width $W$ of the machine's
encoding.
\end{lemma}

\begin{definition}[Increment labels up to a scaling factor]
\lean{Complexity.Codec.scIncs}
\label{Complexity.Codec.scIncs}
\leanok
\uses{Complexity.Codec.Lab, Complexity.Codec.ScId, Complexity.Codec.scFactor, Complexity.Codec.scFactor_le}
For a nondeterministic Turing machine $N$ with Boolean tape alphabet, a scaling
identifier $c$, and a natural number $i$, let $q_c$ be its scaling factor. If $i<q_c$,
assign the label that increments counter $c$ at position $i$; otherwise assign the label
for counter $c$ itself.
\end{definition}

\begin{definition}[Next labels for scaling operations]
\lean{Complexity.Codec.scNext}
\label{Complexity.Codec.scNext}
\leanok
\uses{Complexity.Codec.Lab, Complexity.Codec.ScId, Complexity.Codec.W}
For a nondeterministic Turing machine $N$ with Boolean tape alphabet, define the
next-label map on the scaling identifiers by sending $E$ to the first increment-check
label for counter $E$, $F$ to the first clearing label, $E2$ to the corresponding first
increment-check label for $E2$, $F2$ to the second clearing label, $In$ to the
input-reading label, and $C$ to the first increment-check label for counter $C$.
\end{definition}

\begin{definition}[Registers for increment counters]
\lean{Complexity.Codec.icReg}
\label{Complexity.Codec.icReg}
\leanok
\uses{Complexity.Codec.IcId, Complexity.Codec.rE, Complexity.Codec.rE2, Complexity.Codec.rU}
For a nondeterministic Turing machine $N$ with Boolean tape alphabet, define the
register map on the five increment-counter identifiers $E,U,E2,U2,C$ to take the
respective values $7,14,9,14,14$, among the $21$ registers.
\end{definition}

\begin{definition}[Increment-counter sizes]
\lean{Complexity.Codec.icCount}
\label{Complexity.Codec.icCount}
\leanok
\uses{Complexity.Codec.IcId, Complexity.Codec.W, Complexity.Codec.hH}
For a nondeterministic Turing machine $N$ with Boolean tape alphabet, let $h$ be the
number of pairs consisting of an optional state of $N$ and one of three head positions,
and put $W=h+3k+1$, where $k$ is the number of tapes. Define the sizes for identifiers
$E,U,E2,U2,C$ to be respectively $2W,4,2W,4,2(h+2)$.
\end{definition}

\begin{lemma}[Bound on Increment Counts]
\lean{Complexity.Codec.icCount_le}
\label{Complexity.Codec.icCount_le}
\leanok
\uses{Complexity.Codec.IcId, Complexity.Codec.W, Complexity.Codec.icCount}
\difficulty{4}
Let $N$ be a finite nondeterministic Turing machine with Boolean input alphabet, and let
$k$ be its number of work tapes. Write $H$ for the number of head descriptors and
$W=H+3k+1$. For every increment-count identifier $c$, its associated count is at most
$4W$.
\end{lemma}

\begin{definition}[Next label for an instruction counter]
\lean{Complexity.Codec.icNext}
\label{Complexity.Codec.icNext}
\leanok
\uses{Complexity.Codec.IcId, Complexity.Codec.Lab}
Let $N$ be a finite nondeterministic Turing machine with Boolean tape symbols. The map
from instruction identifiers to labels sends $E$ to the copy label for $F$, $U$ to the
length-check label, $E2$ to the copy label for $F2$, $U2$ to the second length-check
label, and $C$ to the center-read label.
\end{definition}

\begin{definition}[Label at an instruction index]
\lean{Complexity.Codec.icL}
\label{Complexity.Codec.icL}
\leanok
\uses{Complexity.Codec.IcId, Complexity.Codec.Lab, Complexity.Codec.icCount, Complexity.Codec.icCount_le, Complexity.Codec.icNext}
Let $N$ be a finite nondeterministic Turing machine with Boolean tape symbols, let $c$
be an instruction identifier, and let $i$ be a natural number. If $i$ is less than the
instruction count assigned to $c$, the resulting label is the indexed instruction label
at position $i$; otherwise it is the next label assigned to $c$.
\end{definition}

\begin{definition}[Continuation label for a prefix scan]
\lean{Complexity.Codec.prCont}
\label{Complexity.Codec.prCont}
\leanok
\uses{Complexity.Codec.CMode, Complexity.Codec.Lab}
Let $N$ be a finite nondeterministic Turing machine with Boolean tape symbols, let $m$
be one of the four code modes (valid, initial, accepting, or step-pair), let
$p\in\{0,1,2\}$, and let $k\in\{0,1,2,3\}$. For $k=0$ or $1$, the label tests position
$p$ with Boolean parameter false or true, respectively; for $k=2$, it scans at position
$p$. For $k=3$, it selects the copy label for $A$ when $p=0$, the copy label for $B$
when $p=1$ and $m$ is the step-pair mode, the copy label for $Q1s$ when $p=1$ in any
other mode or when $p=2$.
\end{definition}

\begin{definition}[Prefix-scan advancement label]
\lean{Complexity.Codec.prAdvL}
\label{Complexity.Codec.prAdvL}
\leanok
\uses{Complexity.Codec.CMode, Complexity.Codec.Lab, Complexity.Codec.prCont}
Let $N$ be a finite nondeterministic Turing machine with Boolean tape symbols, let $m$
be one of the four code modes (valid, initial, accepting, or step-pair), let
$p\in\{0,1,2\}$ and $k\in\{0,1,2,3\}$, and let $i$ be a natural number. If $i=0$, the
label is the advancement label indexed by $p$ and $k$; if $i>0$, it is the continuation
label determined by $m$, $p$, and $k$.
\end{definition}

\begin{definition}[Second verifier advancement label]
\lean{Complexity.Codec.vA2}
\label{Complexity.Codec.vA2}
\leanok
\uses{Complexity.Codec.Lab, Complexity.Codec.VAd}
Let $N$ be a finite nondeterministic Turing machine with Boolean tape symbols, let $s$
be a verifier address state, and let $i$ be a natural number. If $i<2$, the label is the
second advancement label for $s$, indexed by $i$; otherwise it is the test label for
$s$.
\end{definition}

\begin{definition}[First verifier advancement label]
\lean{Complexity.Codec.vA1}
\label{Complexity.Codec.vA1}
\leanok
\uses{Complexity.Codec.Lab, Complexity.Codec.VAd, Complexity.Codec.vA2}
Let $N$ be a finite nondeterministic Turing machine with Boolean tape symbols, let $s$
be a verifier address state, and let $i$ be a natural number. If $i<2$, the label is the
first advancement label for $s$, indexed by $i$; otherwise it is the second advancement
label for $s$ at index $0$.
\end{definition}

\begin{definition}[Verifier transition function]
\lean{Complexity.Codec.δV}
\leanok
\uses{Complexity.Codec.VAd, Complexity.Codec.adStep, Complexity.Codec.vChunk}
Let $N$ be a finite nondeterministic Turing machine with Boolean tape symbols. The
verifier transition function maps a verifier address state and a pair of Boolean input
symbols to the state obtained by applying one address step using the verifier chunk
associated with $N$.
\end{definition}

\begin{definition}[Simulation transition function]
\lean{Complexity.Codec.δS}
\leanok
\uses{Complexity.Codec.Gs, Complexity.Codec.SAd, Complexity.Codec.SInfo, Complexity.Codec.adStep, Complexity.Codec.sChunk}
Let $N$ be a finite nondeterministic Turing machine with Boolean tape symbols, and let
$I$ consist of a verifier state and an optional Boolean value. The simulation transition
function maps an address state and a pair of Boolean input symbols to the state obtained
by applying one address step using the simulation chunk determined by $N$ and $I$, with
its Boolean parameter set first to false and then to true.
\end{definition}

\begin{definition}[Second simulation advance label]
\lean{Complexity.Codec.sA2}
\label{Complexity.Codec.sA2}
\leanok
\uses{Complexity.Codec.Lab, Complexity.Codec.SAd, Complexity.Codec.SInfo}
For a machine $N$, simulation information $I$, Boolean flags $f_0,f_1$, simulation state
$st$, and natural number $i$, define a label as follows: if $i<2$, it is the second
advance label carrying $I,f_0,f_1,st$ and the index $i$; otherwise it is the simulation
test label carrying $I,f_0,f_1,st$.
\end{definition}

\begin{definition}[First simulation advance label]
\lean{Complexity.Codec.sA1}
\label{Complexity.Codec.sA1}
\leanok
\uses{Complexity.Codec.Lab, Complexity.Codec.SAd, Complexity.Codec.SInfo, Complexity.Codec.sA2}
For a machine $N$, simulation information $I$, Boolean flags $f_0,f_1$, simulation state
$st$, and natural number $i$, define a label as follows: if $i<2$, it is the first
advance label carrying $I,f_0,f_1,st$ and the index $i$; otherwise it is the second
simulation advance label with the same data and index $0$.
\end{definition}

\begin{definition}[Next simulation phase]
\lean{Complexity.Codec.pmNext}
\label{Complexity.Codec.pmNext}
\leanok
\uses{Complexity.Codec.Lab, Complexity.Codec.SInfo, Complexity.Codec.sAd0}
For a machine $N$, simulation information $I$, and Boolean values $b,f,r$, define the
next label to be the simulation test label carrying $I,f,r$ and the initial simulation
state when $b$ is true; when $b$ is false, define it to be the phase label carrying $I$,
flags $\mathrm{true},r$, and starting phase.
\end{definition}

\begin{definition}[Register Program for the Codec]
\lean{Complexity.Codec.CA}
\label{Complexity.Codec.CA}
\leanok
\uses{Complexity.Codec.CMode, Complexity.Codec.Gs, Complexity.Codec.Lab, Complexity.Codec.VAcc, Complexity.Codec.aI, Complexity.Codec.cpDst, Complexity.Codec.cpNext, Complexity.Codec.cpSrc, Complexity.Codec.dcAcc, Complexity.Codec.dcBase, Complexity.Codec.dcDowns, Complexity.Codec.dcIncs, Complexity.Codec.dcNext, Complexity.Codec.dcPtr, Complexity.Codec.goto, Complexity.Codec.icL, Complexity.Codec.icReg, Complexity.Codec.jBit, Complexity.Codec.jLen, Complexity.Codec.pmNext, Complexity.Codec.prAdvL, Complexity.Codec.rB, Complexity.Codec.rCap, Complexity.Codec.rE, Complexity.Codec.rF, Complexity.Codec.rF2, Complexity.Codec.rL, Complexity.Codec.rN1, Complexity.Codec.rP, Complexity.Codec.rPos, Complexity.Codec.rPos2, Complexity.Codec.rQ1, Complexity.Codec.rQ2, Complexity.Codec.rQ3, Complexity.Codec.rQ4, Complexity.Codec.rT, Complexity.Codec.rU, Complexity.Codec.rWp, Complexity.Codec.rZ, Complexity.Codec.sA1, Complexity.Codec.sA2, Complexity.Codec.sDone, Complexity.Codec.scDst, Complexity.Codec.scIncs, Complexity.Codec.scNext, Complexity.Codec.scSrc, Complexity.Codec.vA1, Complexity.Codec.vA2, Complexity.Codec.δS, Complexity.Codec.δV, Complexity.LogProg.ARM}
For a Boolean nondeterministic Turing machine $N$ and a mode $m$ chosen from the four
cases valid, initial, accepting, and step, define a register program with $21$ registers
and two deciders by the instruction assigned to each label. The labels and their
instructions are precisely those in the definition: the program scans and validates
encoded inputs, performs register updates and decider calls, and returns a Boolean
answer according to the selected mode.
\end{definition}
